% !TeX root = ../../relatorio-projeto.tex
% Modelo para preenchimento: substitua as orientacoes por seu relato final.
\chapter{Introdução e contexto}
Apresente o ambiente profissional, a demanda, o problema, a justificativa e
os objetivos do projeto realizado. Delimite o escopo e a participação do
estudante. Use o roteiro em conjunto com as orientações do curso e o manual
\citeind{manualundf}. Identifique as permissões de acesso aos materiais usados.

\chapter{Fundamentação e requisitos}
Articule as referências científicas e técnicas que sustentaram a solução.
Descreva os requisitos, restrições e critérios de aceitação utilizados.
Cadastre suas fontes no arquivo bibliográfico indicado na entrada principal,
substituindo a entrada demonstrativa quando não for citada no trabalho.

\chapter{Metodologia e desenvolvimento realizado}
Relate procedimentos executados, tecnologias, materiais, etapas e decisões.
Explique mudanças de escopo e a contribuição individual. Se apresentar
cronograma, registre o percurso realizado e os desvios relevantes. Escreva
sobre o trabalho concluído, não um plano de execução futura.

\chapter{Solução e entregas}
Descreva o produto, processo ou serviço desenvolvido e seu estágio real de
conclusão. Apresente documentação e artefatos que permitam compreendê-lo.
Identifique os materiais complementares e seus locais de acesso quando cabível.

\chapter{Resultados e avaliação}
Informe critérios, procedimentos, ambiente e resultados da avaliação realizada.
Relacione cada afirmação às evidências e aos objetivos. Distinga entrega,
teste, implantação e impacto observado. Não atribua melhoria operacional
a uma solução que não foi avaliada nesse contexto.

\chapter{Discussão e limitações}
Interprete resultados, decisões e dificuldades. Analise divergências entre
esperado e obtido, resultados negativos e limites de aplicação. Explique o que
as evidências sustentam e quais questões permaneceram abertas.

\chapter{Conclusão}
Responda aos objetivos e sintetize a contribuição técnica e acadêmica.
Apresente limites e recomendações decorrentes da análise. Reserve atividades
futuras às recomendações, sem substituir o relato do que foi realizado.
